% 01_intro
\section{引言与目标}

在边缘设备上运行推理已成常态，模型、调度、驱动、供电各层集中于同一块 SoC（片上系统）之上，操作系统既要为加速器持续供给数据，又要在多核之间分配处理器时间。将这样一台系统从以 C 语言实现、运行于 Linux 的形态迁移到以 Rust 实现的内核，收益在于内存安全与模块化，前提是通用负载的性能不因迁移而下降。所迁移的负载是一台网球拾取机器人，程序移植自 \texttt{pengzechen/\allowbreak aka-rk3588}，原基于 C 语言、运行于 Linux。它的完整回路包含追球、抓取、回桶、寻桶、趋近、投放六个状态，衡量口径为自相机成帧到执行器指令的帧到指令延迟，这也是竞赛所采用的口径。

工作分为两部分。应用层一侧，让机器人在 StarryOS 上运行并\ours{优化其推理链路}，依靠取球循环与数据布局，将帧到指令延迟\keyres{降低至接近 Linux 的水平}。系统层一侧，令整套 StarryOS 作为通用操作系统，在调度、系统调用、内存、频率这些通用的子系统上与 Linux 直接对比。据此，机器人只是众多负载中的一个，所考核的对象是操作系统的通用能力。

比较遵循统一准则。每个子系统采用各自的基准套件，各配一条独立的 Linux 基线，两侧运行同一份遵循 Linux ABI 的可执行文件，一份负载在两侧的差异因此只反映操作系统本身。后文先展开机器人应用，再依次呈现观测、调度、系统调用、内存、频率、启动等系统层子系统。对等类子系统各以对应的 Linux 基线为参照，并叙述相应改动；显示、图形合成、QEMU NPU 模型与板级 I/O 属使能与仿真类贡献，不作对等主张。
